\SourcePage{111}{116}
\section{Spherical waves}

The wave functions of a free spin-$1/2$ particle with definite values of the total angular momentum $j$ are spinor spherical waves. We shall determine their form, first recalling the analogous expressions from non-relativistic theory.

The non-relativistic wave function is a 3-spinor $\psi = \binom{\psi^1}{\psi^2}$. For a state with definite energy $\varepsilon$ (and hence definite magnitude of the momentum $p$\,\footnote{Throughout this section $p$ denotes $|\mathbf{p}|$.}), orbital angular momentum $l$, total angular momentum $j$, and projection $m$, it takes the form
\begin{equation}
\psi = R_{pl}(r)\Omega_{jlm}(\theta, \varphi).
\tag{24.1}
\end{equation}
The angular factor $\Omega_{jlm}$ is a 3-spinor whose components (for the two values $j = l \pm 1/2$ possible at a given $l$) are
\begin{equation}
\Omega_{l+1/2,\,l,\,m} = \begin{pmatrix} \sqrt{\dfrac{j+m}{2j}}\,Y_{l,\,m-1/2}\\[6pt] \sqrt{\dfrac{j-m}{2j}}\,Y_{l,\,m+1/2} \end{pmatrix},
\tag{24.2}
\end{equation}
\[
\Omega_{l-1/2,\,l,\,m} = \begin{pmatrix} -\sqrt{\dfrac{j-m+1}{2j+2}}\,Y_{l,\,m-1/2}\\[6pt] \sqrt{\dfrac{j+m+1}{2j+2}}\,Y_{l,\,m+1/2} \end{pmatrix}
\]
(see~III, \S\,106, problem). We shall call $\Omega_{jlm}$ the \emph{spherical spinors}. Their normalisation is
\begin{equation}
\int \Omega^*_{jlm}\Omega_{j'l'm'}\,do = \delta_{jj'}\delta_{ll'}\delta_{mm'}.
\tag{24.3}
\end{equation}

The radial functions $R_{pl}$ are a common factor in both spinor components of $\psi$ and are given by the formula

\SourcePage{112}{117}
(33.10) (see~III):
\begin{equation}
R_{pl} = \sqrt{\frac{2\pi p}{r}}\,J_{l+1/2}(pr).
\tag{24.4}
\end{equation}
They are normalised according to
\begin{equation}
\int_0^\infty r^2 R_{p'l}R_{pl}\,dr = 2\pi\delta(p' - p).
\tag{24.5}
\end{equation}

Turning now to the relativistic case, we recall first of all that for a moving particle no separate conservation laws for spin and orbital angular momentum exist: the operators $\widehat{\mathbf{s}}$ and $\widehat{\mathbf{l}}$ individually fail to commute with the Hamiltonian. The parity of the state, however, remains conserved (for a free particle). Consequently the quantum number $l$ ceases to specify a definite value of the orbital angular momentum, but it does determine (see below) the parity of the state.

We shall work with the desired wave function (a bispinor) in the standard representation: $\psi = \binom{\varphi}{\chi}$. Under rotations $\varphi$ and $\chi$ transform as 3-spinors, so their angular dependence is again carried by the spherical spinors $\Omega_{jlm}$. Let $\varphi \propto \Omega_{jlm}$, where $l$ is one (definite) of the two values $j + 1/2$ or $j - 1/2$. Under inversion $\varphi(\mathbf{r}) \to i\varphi(-\mathbf{r})$ (see~(21.18)), and since $\Omega_{jlm}(-\mathbf{n}) = (-1)^l\Omega_{jlm}(\mathbf{n})$,
\[
\varphi(\mathbf{r}) \to i(-1)^l\varphi(\mathbf{r}).
\]

The lower component $\chi$, on the other hand, transforms under inversion as $\chi(\mathbf{r}) \to -i\chi(-\mathbf{r})$. For the state to possess a definite parity (i.e.\ for inversion to multiply every component by the same phase), the angular dependence of $\chi$ must involve $\Omega_{jl'm}$ with the other value of $l$: because the two allowed values differ by unity, $(-1)^{l'} = -(-1)^l$.

Furthermore, the radial dependence of $\varphi$ and $\chi$ is governed by the same functions $R_{pl}$ and $R_{pl'}$ (with $l$ and $l'$ matching the order of the spherical harmonics entering $\Omega_{jlm}$). This follows from the fact that each component of $\psi$ satisfies the second-order equation $(\widehat p^2 - m^2)\psi = 0$, which for a given $|\mathbf{p}|$ reads
\[
(\Delta + \mathbf{p}^2)\psi = 0,
\]
coinciding in form with the free-particle Schr\"odinger equation of the non-relativistic theory.

\SourcePage{113}{118}
Thus
\begin{equation}
\varphi = AR_{pl}\Omega_{jlm},\qquad \chi = BR_{pl'}\Omega_{jl'm},
\tag{24.6}
\end{equation}
and it remains to fix the constant coefficients $A$ and $B$. To this end we examine the far zone, where the spherical wave may be treated as a plane wave. By the asymptotic formula~(33.12) (see~III)
\begin{equation}
R_{pl} \approx \frac{1}{ir}\left\{e^{i\bigl(pr - \frac{\pi l}{2}\bigr)} - e^{-i\bigl(pr - \frac{\pi l}{2}\bigr)}\right\},
\tag{24.7}
\end{equation}
so that $\varphi$ is a superposition of two plane waves propagating along $\pm\mathbf{n}$ ($\mathbf{n} = \mathbf{r}/r$). For each of these, by~(23.8),
\[
\chi = \frac{p}{\varepsilon + m}(\pm\mathbf{n}\boldsymbol{\sigma})\varphi.
\]

From what has been said above (formulae~(24.6)) it is clear that $(\mathbf{n}\boldsymbol{\sigma})\Omega_{jlm} = a\Omega_{jl'm}$ with $a$ a constant. This constant is easily found by evaluating both sides at $m = 1/2$ with $\mathbf{n}$ along the $z$-axis. Using~(7.2a) one obtains
\begin{equation}
(\mathbf{n}\boldsymbol{\sigma})\Omega_{jlm} = i^{l'-l}\Omega_{jl'm}.
\tag{24.8}
\end{equation}

Collecting the formulae and comparing with~(24.6), we find
\[
B = -\frac{p}{\varepsilon + m}\,A.
\]
Finally, the coefficient $A$ is fixed by the overall normalisation of $\psi$. Adopting the condition
\begin{equation}
\int \psi^*_{pjlm}\psi_{p'j'l'm'}\,d^3x = 2\pi\delta_{jj'}\delta_{ll'}\delta_{mm'}\delta(p - p'),
\tag{24.9}
\end{equation}
we arrive at the final result
\begin{equation}
\psi_{pjlm} = \frac{1}{\sqrt{2\varepsilon}}\begin{pmatrix} \sqrt{\varepsilon + m}\,R_{pl}\,\Omega_{jlm}\\ -\sqrt{\varepsilon - m}\,R_{pl'}\,\Omega_{jl'm} \end{pmatrix},\qquad l' = 2j - l.
\tag{24.10}
\end{equation}

At fixed $j$ and $m$ (and energy $\varepsilon$) there are therefore two states that differ in parity. Which parity a state carries is uniquely determined by $l$ ($= j \pm 1/2$): inversion multiplies the bispinor~(24.10) by $i(-1)^l$. The upper and lower components of this bispinor do, however, involve spherical harmonics of both orders $l$ and $l'$---a direct manifestation of the fact that the orbital angular momentum is not separately well-defined.

As $r \to \infty$ one can, in any small region of space, regard the spherical waves~(24.7) as plane waves with momentum $\mathbf{p} = \pm p\mathbf{n}$. It is therefore clear that the momentum-space wave functions differ from~(24.10) mainly in the

\SourcePage{114}{119}
absence of radial factors, with $\mathbf{n}$ now denoting the direction of the momentum.

To pass directly to the momentum representation one performs the Fourier transform:
\begin{equation}
\psi(\mathbf{p}') = \int \psi(\mathbf{r})\,e^{-i\mathbf{p}'\mathbf{r}}\,d^3x,
\tag{24.11}
\end{equation}

The integral is evaluated with the aid of the plane-wave expansion in spherical harmonics (see~III, (34.3)):
\begin{equation}
e^{i\mathbf{p}\mathbf{r}} = \frac{2\pi}{p}\sum_{l=0}^{\infty}\sum_{m=-l}^{l}i^l R_{pl}(r)\,Y^*_{lm}\Bigl(\frac{\mathbf{p}}{p}\Bigr)Y_{lm}\Bigl(\frac{\mathbf{r}}{r}\Bigr).
\tag{24.12}
\end{equation}
Substituting $e^{-i\mathbf{p}'\mathbf{r}}$ into~(24.11) via this expansion and using~(24.5), the Fourier components of the function
\[
\psi(\mathbf{r}) = R_{pl}(r)\,\Omega_{jlm}\Bigl(\frac{\mathbf{r}}{r}\Bigr)
\]
are found to be
\[
\psi(\mathbf{p}') = \frac{(2\pi)^2}{p}\,\delta(p' - p)\,i^{-l}Y_{lm'}\Bigl(\frac{\mathbf{p}'}{p'}\Bigr)\int \Omega_{jlm}\Bigl(\frac{\mathbf{r}}{r}\Bigr)Y^*_{lm'}\Bigl(\frac{\mathbf{r}}{r}\Bigr)\,do.
\]
The integral here equals the coefficients of the spherical harmonics in the definition of the spherical spinors~(24.2), and together with the factor $Y_{lm'}(\mathbf{p}'/p')$ it reconstructs the same spherical spinor, now depending on $\mathbf{p}'/p'$:
\[
\psi(\mathbf{p}') = \frac{(2\pi)^2}{p}\,\delta(p' - p)\,i^{-l}\,\Omega_{jlm}\Bigl(\frac{\mathbf{p}'}{p'}\Bigr).
\]

Applying this result to the bispinor wave function~(24.10), we obtain its momentum-space form
\begin{equation}
\psi_{pjlm}(\mathbf{p}') = \delta(p' - p)\,\frac{(2\pi)^2}{p\sqrt{2\varepsilon}}\begin{pmatrix} \sqrt{\varepsilon + m}\,i^{-l}\,\Omega_{jlm}(\mathbf{p}'/p')\\ -\sqrt{\varepsilon - m}\,i^{-l'}\,\Omega_{jl'm}(\mathbf{p}'/p') \end{pmatrix}.
\tag{24.13}
\end{equation}

The states $|pjlm\rangle$ coincide with the states $|pjm|\lambda||$ considered in \S\,16 (with $|\lambda| = 1/2$): both sets possess definite values of $pjm$ and parity. The spherical spinors $\Omega_{jlm}$ are therefore expressible through the functions $D^{(j)}_m$ (both depending on the argument $\mathbf{p}/p$). In the limit $p \to 0$ the wave functions~(24.13) reduce to 3-spinors $\Omega_{jlm}$ with parity $P = \eta(-1)^l$ (where $\eta = i$ is the ``intrinsic parity'' of the spinor). Comparison with the results of \S\,16 leads to the formula
\begin{equation}
\Omega_{jlm} = i^l\sqrt{\frac{2j+1}{8\pi}}\Bigl(w^{(-1/2)}D^{(j)}_{-1/2,m} \pm w^{(1/2)}D^{(j)}_{1/2,m}\Bigr)
\tag{24.14}
\end{equation}
(with $l = j \mp 1/2$, where $w^{(\lambda)}$ are the 3-spinors~(23.14)).
